\begin{tikzpicture}[
    thick,
    % Stile per i fermioni (linea continua con freccia a metà)
    fermion/.style={draw=black, postaction={decorate},
        decoration={markings, mark=at position 0.55 with {\arrow{latex}}}},
    % Stile per il campo scalare (linea tratteggiata)
    scalar/.style={draw=black, dashed},
    % Stile per la X
    cross/.style={cross out, draw=black, thick, minimum size=8pt, inner sep=0pt},
    % Stile per il vertice puntiforme
    dot/.style={circle, fill=black, minimum size=4pt, inner sep=0pt}
]

% ---------------- DIAGRAMMA DI DIRAC (Sinistra) ----------------
\coordinate (D) at (0,0);

% Linee entrante e uscente
\draw[fermion] (-1.5, 1.5) node[left] {$\nu_L$} -- (D);
\draw[fermion] (D) -- (1.5, 1.5) node[right] {$\nu_R$};

% Linea tratteggiata (Higgs) con la X finale
\draw[scalar] (D) -- (0, -1.5);
\node[cross] at (0, -1.5) {};

% Vertice centrale con pallino e label
\node[dot, label=right:{$m_D$}] at (D) {};


% ---------------- DIAGRAMMA DI MAJORANA (Destra) ----------------
\coordinate (M) at (6, 0.5); % Spostato a destra e leggermente in alto per allineamento visivo

% Linee (entrambe puntano verso il centro)
\draw[fermion] (4, 0.5) node[left] {$\nu_R$} -- (M);
\draw[fermion] (8, 0.5) node[right] {$\overline{\nu}_L$} -- (M);

% Vertice centrale con la X e label
\node[cross, label=above:{$M$}] at (M) {};

\end{tikzpicture}